\chapter{Description of the Study Area}
\label{wasenmoos}

The study area, the nature reserve (German: Naturschutzgebiet) - Wasenmoos is a Bog Wetland located in Tettnang, Baden-Württemberg, Germany. Exhaustive information on the nature reserve (in German) may be found here \url{https://rips-dienste.lubw.baden-wuerttemberg.de/rips/ripsservices/apps/naturschutz/schutzgebiete/steckbrief.aspx?id=939001000068}.
\\ \\
The demarcated study area as seen in Figures \ref{fig:wasenmoos_map_closeup} and \ref{fig:wasenmoos_qgis} lies in a classified Bog type wetland area. In short, these are areas formed by historical peat deposits consisting of decaying and compressed biomass mostly produced by the vegetation it has hosted in the past. The area is home to various wetland vegetative species that in turn support corresponding animal life.
\\ \\
Over the years, the bog wetland has reduced in area resulting from their dewatering via strategically constructed drainage so that they may be used for farming. An ongoing study aims at studying the effects of these streams on the vegetation in the nature reserve by numerically modelling the study area and estimating the soil moisture. This study aims to apply the conceptualized model to the study area and aid the former study in coupling the GW and SW models and in estimating the soil moisture.
\\ \\

\begin{figure}[h!]
    \centering
    \includegraphics[width=\linewidth]{\Pics/illustrations/wasenmoos_map}
    \caption{Location of the study area within southern Germany.}
    \label{fig:wasenmoos_map}
\end{figure}


\begin{figure}[h!]
    \centering
    \includegraphics[width=\linewidth]{\Pics/illustrations/wasenmoos_map_closeup}
    \caption{Map of the study area.}
    \label{fig:wasenmoos_map_closeup}
\end{figure}

\begin{figure}[h!]
    \centering
    \includegraphics[width=\linewidth]{\Pics/illustrations/wasenmoos_qgis}
    \caption{Topological Map showing the demarcated study area (in red), the streams (in blue), and the GW observation locations (in green).}
    \label{fig:wasenmoos_qgis}
\end{figure}

As is observed from Figure \ref{fig:wasenmoos_qgis}, the study area is relatively flat and bordered by ridges on both sides, although the ridges on the right edge are more pronounced. Figure \ref{fig:wasenmoosCS} illustrates the hydrogeological aspects of a typical cross-section of the study area. Here, the unconfined aquifer, composed mainly of peat inside the study area, is roughly estimated to be \unit[0.5]{m} to \unit[3.5]{m} thick. Following which lies an aquitard composed of clay, estimated to be roughly around \unit[1]{m} thick. Underneath which lies a confined Quaternary aquifer composed of fluvial sediments with an estimated thickness of somewhere between \unit[10]{m} to \unit[15]{m}. The dewatering stream/drain that touches the study area is believed to be dug to the depth of the start of the aquitard.
\\
\begin{figure}[h!]
    \centering
    \includegraphics[width=\linewidth]{\Pics/illustrations/wasenmoosCS}
    \caption{Hydrogeological representation: A typical cross-section of the study area; not to scale.}
    \label{fig:wasenmoosCS}
\end{figure}


\endinput